% Template for post or presentation abstracts submitted to ILAS 2022
% 1. Fill in the presenter's information (name, affiliation, email
% address) and talk information (title of presentation, abstract).
% 2. Compile to make sure there are no errors.
% 3. Save the .tex file for your abstract with a distinctive name,
% ideally "FamilynameGivenname.tex".
% 4. Upload your .tex file to https://clr.ie/129557
%       
% * Please keep your LaTeX commands simple, and avoid the use of
%    user-defined macros.
% * Include details of co-authors and funding acknowledgements after the abstract.
% * Upload only the LaTeX file (with .tex extension).
% * Questions: Contact Niall Madden (Niall.Madden@NUIGalway.ie)

\documentclass[a4paper]{amsart} 
\usepackage{amssymb} 
\usepackage{hyperref}
\pagestyle{empty}
\date{} 

%% IGNORE THE NEXT 4 LINES
\makeatletter % 
\def\labelenumi{\theenumi.}
\def\theenumi{\@arabic\c@enumi}
\makeatother
%% STOP IGNORING NOW

\begin{document}

\title{Comparison Theorems for Splittings of M-matrices in block Hessenberg Form}
\author{Luca Gemignani} %Presenting author only
\address{University of Pisa} % Affiliation of presenting author
\email{luca.gemignani@unipi.it} % Email address of presenting author only

\maketitle

% The abstract  ***no more than one page***

In this  talk  we consider the solution of $M$-matrix linear systems  
 in block Hessenberg form, and we show new comparison results among 
 matrix splittings that hold for this special structure.
In particular, we prove that for
a lower-Hessenberg M-matrix $\rho(P_{GS})\geq \rho(P_S)\geq \rho(P_{AGS})$,
 where $\rho(A)$ denotes the spectral radius of $A$ and 
$P_{GS}, P_S, P_{AGS}$ are the iteration matrices of the Gauss--Seidel,
staircase, and anti-Gauss--Seidel method.
This is a result that does not seem to follow from classical
comparison results, as these splittings are not directly comparable.
Also,  it  fosters the  use  of stair partitionings  for solving  Hessenberg  linear systems in parallel.  


\emph{This is joint work with Federico Poloni (Pisa).}
 
\end{document}


% Please leave the next bit alone  
%%% Local Variables: 
%%% mode: latex
%%% TeX-master: "../ILAS-2022"
%%% End:
